% =====================================================================
%  Список обозначений. Добавляя новый макрос в preamble/macros.tex,
%  не забудьте дописать его сюда.
% =====================================================================
\section*{Основные обозначения}
\addcontentsline{toc}{section}{Основные обозначения}

\newcommand{\notationgroup}[1]{%
  \midrule\multicolumn{2}{@{}l}{\textbf{#1}}\\[0.2em]}

\begingroup
\small
\renewcommand{\arraystretch}{1.1}
\begin{longtable}{@{} >{$\displaystyle}p{0.28\textwidth}<{$}
                      >{\raggedright\arraybackslash}p{0.68\textwidth} @{}}
  \toprule
  \multicolumn{1}{@{}l}{\textbf{Обозначение}} & \textbf{Смысл} \\
  \endfirsthead
  \toprule
  \multicolumn{1}{@{}l}{\textbf{Обозначение}} & \textbf{Смысл} \\
  \endhead
  \bottomrule
  \endfoot

  \notationgroup{Множества}
  \R,\ \N,\ \Z            & вещественные, натуральные и целые числа \\
  \mathcal{B}(\R^n)       & борелевская $\sigma$-алгебра на $\R^n$ \\

  \notationgroup{Вероятность и моменты}
  \P(A)                   & вероятность события $A$ \\
  \E X                    & математическое ожидание \\
  \D X = \Var X           & дисперсия: $\E (X - \E X)^2$ \\
  \cov(X, Y)              & ковариация: $\E (X - \E X)(Y - \E Y)$ \\
  \corr(X, Y)             & коэффициент корреляции: $\frac{\cov(X, Y)}{\sqrt{\D X \, \D Y}}$ \\
  \Ind{A}                 & индикатор события $A$ (равен $1$, если $A$ произошло, иначе $0$) \\

  \notationgroup{Выборка и статистическая модель}
  X_{[n]} = (\sample)     & выборка размера $n$ \\
  X_i \iid F              & $X_i$ независимы и одинаково распределены с функцией распределения $F$ \\
  \Th                     & параметрическое множество \\
  \{\Pt : \theta \in \Th\} & параметрическое семейство распределений \\
  \Pt,\ \Et,\ \Dt         & вероятность, матожидание и дисперсия при значении параметра $\theta$ \\
  \est\theta              & оценка параметра $\theta$ (функция от выборки) \\
  \Xbar                   & выборочное среднее: $\frac1n \sum_{i=1}^n X_i$ \\
  a_k^*,\ a_k^{0*}        & начальный и центральный эмпирические моменты порядка $k$ \\
  s^2 = a_2^{0*}          & выборочная дисперсия: $\frac1n \sum_{i=1}^n (X_i - \Xbar)^2$ \\
  \tilde{s}^2             & исправленная дисперсия: $\frac{n}{n-1}\, s^2$ \\
  x_p                     & выборочный квантиль порядка $p$; $Q_1, Q_2, Q_3$~--- квартили \\
  X_{(1)} \le \ldots \le X_{(n)} & вариационный ряд ($X_{(k)}$~--- $k$-я порядковая статистика) \\
  \Pemp(B)                & эмпирическое распределение: доля элементов выборки, попавших в $B$ \\
  F_n^*(x)                & эмпирическая функция распределения: $\Pemp\bigl((-\infty, x]\bigr)$ \\
  f_n^*(x)                & гистограмма \\

  \notationgroup{Сходимости случайных величин}
  \xi_n \toas \xi         & сходимость почти наверное: $\P(\xi_n \to \xi) = 1$ \\
  \xi_n \toP \xi          & сходимость по вероятности: $\forall \epsilon > 0 \;\; \P(\abs{\xi_n - \xi} > \epsilon) \to 0$ \\
  \xi_n \toLp{p} \xi      & сходимость в $L_p$: $\E \abs{\xi_n - \xi}^p \to 0$ \\
  \xi_n \tod \xi          & сходимость по распределению: $F_{\xi_n}(x) \to F_\xi(x)$ во всех точках непрерывности $F_\xi$ \\

  \notationgroup{Распределения}
  \Norm(\mu, \sigma^2)    & нормальное с параметрами $\mu$ (среднее) и $\sigma^2$ (дисперсия) \\
  \Bern(p)                & Бернулли: $\P(X = 1) = p$, $\P(X = 0) = 1 - p$ \\
  \Bin(n, p)              & биномиальное: число успехов в $n$ испытаниях Бернулли \\
  \Pois(\lambda)          & пуассоновское: $\P(X = k) = e^{-\lambda} \frac{\lambda^k}{k!}$, $k = 0, 1, \ldots$ \\
  \Exp(\lambda)           & экспоненциальное с плотностью $\lambda e^{-\lambda x}$, $x \ge 0$ \\
  \Unif[a, b]             & равномерное на отрезке $[a, b]$ \\

  \notationgroup{Прочее}
  a \eqdef b              & $a$ по определению равно $b$ \\
  \argmax_\theta f(\theta) & точка, в которой $f$ достигает максимума \\
\end{longtable}
\endgroup

\clearpage
